La miopia és una alteració ocular en la qual la persona veu borrosos els objectes llunyans. Un ull miop és massa convergent i crea la imatge d'objectes llunyans per davant de la retina. La miopia de l'ull es pot rectificar amb una lent correctora. El Joan és miop i sols pot enfocar bé objectes situats a una distància màxima de 4 metres.

\begin{apartat}{1,25}
Quin tipus de lent és necessària per corregir la miopia? Calcula la focal i la potència de la lent correctora per a que el Joan pugui veure focalitzat un objecte llunyà (a l'infinit).
\end{apartat}

En Joan veu un arbre situat a 12 metres d'ell mentre porta posades les seves ulleres

\begin{apartat}{1,25}
Troba a quina distància d'en Joan es forma la imatge de l'arbre creada per la lent correctora. Fes el diagrama de rajos a escala en la següent quadrícula per localitzar gràficament la imatge de l'arbre.
\end{apartat}

\figura[0.95]{fig1.png}
